\documentclass{article}
\usepackage{graphicx}
\usepackage{longtable}
\usepackage{hyperref}
\title{Automated Linux System Provisioning Progress Tracking}
\author{Your Name}
\date{\today}

\begin{document}

\maketitle

\section*{Overview}
This document tracks the progress of the Automated Linux System Provisioning project. Each step is divided into tasks, and detailed descriptions, images, logs, and findings can be documented here.

\section*{Progress Tracking}

\begin{longtable}{|p{4cm}|p{10cm}|}
\hline
\textbf{Step} & \textbf{Description and Progress} \\
\hline
Step 1: PXE Boot Environment Setup & \\ \hline
TFTP Server Setup & Description, issues encountered, solutions, etc. \\ \hline
Initially, Arch Linux doesn't have TFTPD-HPA package, so I had to go searching for it.
I eventually found it under  tftp-hpa not tftpd.
I also needed to make sure inetutils, dnsmasq, and syslinux were installed for DHCP and pxelinux.0
DHCP Server Setup & Description, issues encountered, solutions, etc. \\ \hline
Bootloader Configuration & Description, issues encountered, solutions, etc. \\ \hline
Testing PXE Boot & Description, issues encountered, solutions, etc. \\ \hline
Step 2: Automating OS Installation & \\ \hline
Rocky Linux Installation Automation & Description, issues encountered, solutions, etc. \\ \hline
Arch Linux Installation Automation & Description, issues encountered, solutions, etc. \\ \hline
Step 3: Post-Install Configuration with Ansible & Description, issues encountered, solutions, etc. \\ \hline
Step 4: Management Interface & Description, issues encountered, solutions, etc. \\ \hline
Step 5: Deployment & Documentation & Description, issues encountered, solutions, etc. \\ \hline
\end{longtable}

\end{document}
